\section*{Problem 1}

将
\[
Q\left(\boldsymbol{x}\right) = x_1^2 + 2x_2^2 + 3x_3^2 + 2x_1x_2 + 4x_1x_3 + 2x_2x_3
\]
化为标准形.

\begin{proof}
直接配方得
\[
\begin{aligned}
    Q\left(\boldsymbol{x}\right) &= x_1^2 + 2x_2^2 + 3x_3^2 + 2x_1x_2 + 4x_1x_3 + 2x_2x_3, \\
    &= \left(x_1 + x_2 + 2x_3\right)^2 + x_2^2 - x_3^2 - 2x_2x_3, \\
    &= \boxed{\left(x_1 + x_2 + 2x_3\right)^2 + \left(x_2 - x_3\right)^2 - 2x_3^2}. \\
\end{aligned}
\]

\end{proof}

\section*{Problem 2}

设
\[
Q\left(x_1, x_2, x_3\right) = ax_1^2 + bx_2^2 + ax_3^2 + 2cx_1x_3,
\]
求 $a, b, c$ 满足什么条件时 $Q$ 为正定二次型.

\begin{proof}
若 $Q$ 正定，则 $a=Q(1,0,0)>0$、$b=Q(0,1,0)>0$。在 $a>0$ 下配方得
\[
\begin{aligned}
    Q \left(x_1, x_2, x_3\right) &= a \left(x_1 + \frac{cx_3}{a}\right)^2 + \left(a - \frac{c^2}{a}\right)x_3^2 + bx_2^2, \\
\end{aligned}
\]
因此 $Q$ 正定当且仅当三个平方项的系数均为正，即
\[
\boxed{
    a > 0,\quad b > 0,\quad a^2 - c^2 > 0
}.
\]
等价地，条件可写为 $\boxed{b > 0,\ a > |c|}$。

\end{proof}
\section*{Problem 3}

设
\[
Q \left(\boldsymbol{x}\right) = x_1^2 + 2x_2^2 + 3x_3^2 - 2x_1x_2 + 4x_2x_3
\]
求其正负惯性指数.

\begin{proof}
显然有
\[
\begin{aligned}
    Q \left(\boldsymbol{x}\right) &= x_1^2 + 2x_2^2 + 3x_3^2 - 2x_1x_2 + 4x_2x_3, \\
    &= \left(x_1 - x_2\right)^2 + \left(x_2 + 2x_3\right)^2 - x_3^2. \\
\end{aligned}
\]
因此 $Q$ 的正惯性指数为 $2$, 负惯性指数为 $1$. 即
\[
\boxed{p = 2, q = 1}.
\]

\end{proof}
\section*{Problem 4}

设
\[
A =
\begin{bmatrix}
    2 & -1 & 1 \\
    -1 & 2 & -1 \\
    1 & -1 & 0 \\
\end{bmatrix}
\]
\begin{enumerate}
    \item 构造 $P$ 使 $P^\top A P$ 为标准形；
    \item 求惯性指数；
    \item 判断是否正定.
\end{enumerate}

\begin{proof}
考虑二次型
\[
\begin{aligned}
    Q \left(\boldsymbol{x}\right) &= \boldsymbol{x}^\top A \boldsymbol{x} = \left(x_1, x_2, x_3\right) \begin{bmatrix}
        2 & -1 & 1 \\
        -1 & 2 & -1 \\
        1 & -1 & 0 \\
    \end{bmatrix} \left(x_1, x_2, x_3\right)^\top, \\
    &= 2x_1^2 + 2x_2^2 - 2x_1x_2 - 2x_2x_3 + 2x_1x_3, \\
    &= 2 \left(x_1 - \frac{x_2}{2} + \frac{x_3}{2}\right)^2 + \frac{3x_2^2}{2} - 2x_2x_3 - \frac{x_3^2}{2}, \\
    &= 2 \left(x_1 - \frac{x_2}{2} + \frac{x_3}{2}\right)^2 + \frac{3}{2} \left(x_2 - \frac{x_3}{3}\right)^2 - \frac{2x_3^2}{3}, \\
    &= \left(x_1 - \frac{x_2}{2} + \frac{x_3}{2}, x_2 - \frac{x_3}{3}, x_3\right)
    \mathrm{diag}\left(2, \frac{3}{2}, -\frac{2}{3}\right) \left(x_1 - \frac{x_2}{2} + \frac{x_3}{2}, x_2 - \frac{x_3}{3}, x_3\right)^\top, \\
    &= \left[\begin{bmatrix}
        1 & -\frac{1}{2} & \frac{1}{2} \\
        0 & 1 & -\frac{1}{3} \\
        0 & 0 & 1 \\
    \end{bmatrix} \left(x_1, x_2, x_3\right)^\top\right]^\top \mathrm{diag}\left(2, \frac{3}{2}, -\frac{2}{3}\right) \left[\begin{bmatrix}
        1 & -\frac{1}{2} & \frac{1}{2} \\
        0 & 1 & -\frac{1}{3} \\
        0 & 0 & 1 \\
    \end{bmatrix} \left(x_1, x_2, x_3\right)^\top\right], \\
    &= \left(x_1, x_2, x_3\right) \underbrace{\begin{bmatrix}
        1 & -\frac{1}{2} & \frac{1}{2} \\
        0 & 1 & -\frac{1}{3} \\
        0 & 0 & 1 \\
    \end{bmatrix}^\top}_{{P^{-1}}^\top} \mathrm{diag}\left(2, \frac{3}{2}, -\frac{2}{3}\right) \underbrace{\begin{bmatrix}
        1 & -\frac{1}{2} & \frac{1}{2} \\
        0 & 1 & -\frac{1}{3} \\
        0 & 0 & 1 \\
    \end{bmatrix}}_{P^{-1}} \left(x_1, x_2, x_3\right)^\top.
\end{aligned}
\]
则 $P = \boxed{\begin{bmatrix}
    1 & \frac{1}{2} & -\frac{1}{3} \\
    0 & 1 & \frac{1}{3} \\
    0 & 0 & 1 \\
\end{bmatrix}}$，观察对角阵 $\mathrm{diag}\left(2, \frac{3}{2}, -\frac{2}{3}\right)$ 得 $\boxed{p = 2, q = 1}$，且 $\boxed{A \textrm{ 不是正定矩阵}}$.

\end{proof}
